\section{AIME I 1991}
        \begin{enumerate}
        	\item (\textit{AIME I 1991}) Find $x^2+y^2_{}$ if $x_{}^{}$ and $y_{}^{}$ are positive integers such that

 \begin{align*} xy+x+y&=71, \\ x^2y+xy^2&=880. \end{align*}



 
 \textbf{Answer: }146



\textbf{Solution 1}

Define $a = x + y$ and $b = xy$. Then $a + b = 71$ and $ab = 880$. Solving these two equations yields a \href{https://artofproblemsolving.com/wiki/index.php?title=Quadratic\_equation}{quadratic}: $a^2 - 71a + 880 = 0$, which \href{https://artofproblemsolving.com/wiki/index.php?title=Factor}{factors} to $(a - 16)(a - 55) = 0$. Either $a = 16$ and $b = 55$ or $a = 55$ and $b = 16$. For the first case, it is easy to see that $(x,y)$ can be $(5,11)$ (or vice versa). In the second case, since all factors of $16$ must be $\le 16$, no two factors of $16$ can sum greater than $32$, and so there are no integral solutions for $(x,y)$. The solution is $5^2 + 11^2 = \boxed{146}$.

 





\textbf{Solution 2}

Since $xy + x + y + 1 = 72$, this can be factored to $(x + 1)(y + 1) = 72$. As $x$ and $y$ are \href{https://artofproblemsolving.com/wiki/index.php?title=Integer}{integers}, the possible sets for $(x,y)$ (ignoring cases where $x > y$ since it is symmetrical) are $(1, 35),\ (2, 23),\ (3, 17),\ (5, 11),\ (7,8)$. The second equation factors to $(x + y)xy = 880 = 2^4 \cdot 5 \cdot 11$. The only set with a factor of $11$ is $(5,11)$, and checking shows that it is correct.

 





\textbf{Solution 3}

Let $a=x+y$, $b=xy$ then we get the equations
 \begin{align*} a+b&=71\\ ab&=880 \end{align*}
After finding the \href{https://artofproblemsolving.com/wiki/index.php?title=Prime\_factorization}{prime factorization} of $880=2^4\cdot5\cdot11$, it's easy to obtain the solution $(a,b)=(16,55)$. Thus
 \[x^2+y^2=(x+y)^2-2xy=a^2-2b=16^2-2\cdot55=\boxed{146}\]Note that if $(a,b)=(55,16)$, the answer would exceed $999$ which is invalid for an AIME answer.~ Nafer

 





\textbf{Solution 4}

From the first equation, we know $x+y=71-xy$. We factor the second equation as $xy(71-xy)=880$. Let $a=xy$ and rearranging we get $a^2-71a+880=(a-16)(a-55)=0$. We have two cases: (1) $x+y=16$ and $xy=55$ OR (2) $x+y=55$ and $xy=16$. We find the former is true for $(x,y) = (5,11)$. $x^2+y^2=121+25=146$.

 





\textbf{Solution 5}

First, notice that you can factor $x^2y + xy^2$ as $xy(x + y)$. From this, we notice that $xy$ and $x + y$ is a common occurrence, so that lends itself to a simple solution by substitution. Let $xy = b$ and $x + y = a$. From this substitution, we get the following system:
 \[a + b = 71\]
 \[ab = 880\]Solving that system gives us the following two pairs $(a, b)$: $(16, 55)$ and $(55, 16)$. The second one is obviously too big as $55^2$ is clearly out of bounds for an AIME problem (More on that later). Therefore, we proceed with substituting back the pair $(16, 55)$. This means that $x + y = 16$ and $xy = 55$Then, instead of solving the system, we can do a clever manipulation by squaring $x + y$. Doing so, we get:
 \[(x + y)^2 = (x^2 + y^2) + 2xy\]We see that in this form, we can substitute everything in except for $(x^2 + y^2)$, which is the desired answer. Substituting, we get:
 \[256 = (x^2 + y^2) + 110\]so $x^2 + y^2 = \boxed{146}$. (If we were to go with the pair $(55, 16)$, then the $(x + y)^2$ would be absurdly out of bounds)

 




	\item (\textit{AIME I 1991}) Rectangle $ABCD_{}^{}$ has sides $\overline {AB}$ of length 4 and $\overline {CB}$ of length 3. Divide $\overline {AB}$ into 168 congruent segments with points $A_{}^{}=P_0, P_1, \ldots, P_{168}=B$, and divide $\overline {CB}$ into 168 congruent segments with points $C_{}^{}=Q_0, Q_1, \ldots, Q_{168}=B$. For $1_{}^{} \le k \le 167$, draw the segments $\overline {P_kQ_k}$. Repeat this construction on the sides $\overline {AD}$ and $\overline {CD}$, and then draw the diagonal $\overline {AC}$. Find the sum of the lengths of the 335 parallel segments drawn.


 
 \textbf{Answer: }840



\textbf{Solution 1}


                    \begin{figure}[H]
                        \centering
                        \includegraphics[width=0.5\textwidth]{Solution/AIME_1991_I_2_Solution_Figure_1.png}
                    \end{figure}
                    The length of the diagonal is $\sqrt{3^2 + 4^2} = 5$ (a 3-4-5 \href{https://artofproblemsolving.com/wiki/index.php?title=Right\_triangle}{right triangle}). For each $k$, $\overline{P_kQ_k}$ is the \href{https://artofproblemsolving.com/wiki/index.php?title=Hypotenuse}{hypotenuse} of a $3-4-5$ right triangle with sides of $3 \cdot \frac{168-k}{168}, 4 \cdot \frac{168-k}{168}$. Thus, its length is $5 \cdot \frac{168-k}{168}$. Let $a_k=\frac{5(168-k)}{168}$. We want to find $2\sum\limits_{k=1}^{168} a_k-5$ since we are over counting the diagonal.$2\sum\limits_{k=1}^{168} \frac{5(168-k)}{168}-5 =2\frac{(0+5)\cdot169}{2}-5 =168\cdot5 =\boxed{840}$

 





\textbf{Solution 2}

Using the above diagram, we have that $\Delta ABC \sim \Delta P_k B Q_k$ and each one of these is a dilated 3-4-5 right triangle (This is true since $\Delta ABC$ is a 3-4-5 right triangle). Now, for all $k$, we have that $\overline{P_k Q_k}$ is the hypotenuse for the triangle $P_k B Q_k$. Therefore we want to know the sum of the lengths of all $\overline{P_k Q_k}$.This is given by the following:
 \[2 \cdot(\sum_{k=1}^{168} P_kQ_k) + 5\]

 
 \[= 2 \cdot \frac{ 0+5+10+...+835}{168} +5\]Then by the summation formula for the sum of the terms of an arithmetic series,
 \[= \frac{835 \cdot 168}{168} +5 = 835+5 = \boxed{840}\]

 





\textbf{Solution 3}

First, count the diagonal which has length $5$. For the rest of the segments, think about pairing them up so that each pair makes $5$. For example, the parallel lines closest to the diagonal would have length $\frac{167}{168}\cdot{5}$ while the parallel line closest to the corner of the rectangle would have length $\frac{1}{168}\cdot{5}$ by similar triangles. If you add the two lengths together, it is $\frac{167}{168}\cdot{5} + \frac{1}{168}\cdot{5} = 5.$ There are $\frac{335-1}{2}$ pairs of these segments, for a total of $5+(167)(5)=168(5)=\boxed{840}.$~justlearningmathog

 




	\item (\textit{AIME I 1991}) Expanding $(1+0.2)^{1000}_{}$ by the binomial theorem and doing no further manipulation gives


 ${1000 \choose 0}(0.2)^0+{1000 \choose 1}(0.2)^1+{1000 \choose 2}(0.2)^2+\cdots+{1000 \choose 1000}(0.2)^{1000}$$= A_0 + A_1 + A_2 + \cdots + A_{1000},$where $A_k = {1000 \choose k}(0.2)^k$ for $k = 0,1,2,\ldots,1000$. For which $k_{}^{}$ is $A_k^{}$ the largest?


 
 \textbf{Answer: }166



\textbf{Solution 1 (Proof Format)}

$\textbf{Lemma:}$

 The "maxima" is the point \( k \) in which every \( n > k \) results in a decrease of value.

 $\textbf{Conjecture:}$

 For all values \( k \) upon \( A_k \) there exists a value \( n=k \) in which 

 
 \[\binom{1000}{n} \cdot \left(\frac{1}{5}\right)^n\]

 is as great as possible.

 $\textbf{Proof of the Conjecture}$

 For simplicity, denote $a_k = \binom{1000}{k} \cdot \left(\frac{1}{5}\right)^k$.

 Examine the ratio $\frac{a_{k+1}}{a_k}$. We see this simplifies into 

 
 \[\frac{a_{k+1}}{a_k} =  \frac{\binom{1000}{k+1}\left(\frac{1}{5}\right)^{k+1}}{\binom{1000}{k}\left(\frac{1}{5}\right)^k} = \frac{\binom{1000}{k+1}}{\binom{1000}{k}} \cdot \frac{1}{5}\]

 Then, the identity $\frac{\binom{n}{k+1}}{\binom{n}{k}} = \frac{n-k}{k+1}$ simplifies the above ratio as

 
 \[\frac{a_{k+1}}{a_k} = \frac{1000-k}{k+1} \cdot \frac{1}{5}\]

 which can be rewritten as

 
 \[\frac{1000-k}{5(k+1)}\]

 Note that \( 5(k+1) \) will exponentially increase greater than \( 1000-k \) through properties of multiplication.

 So, for the sequence to begin descent, we want the exact value of \( k \) in which

 \( 1000-k \le 5(k+1) \)

 Solving for which we obtain

 \( k \ge 165 \frac{5}{6} \). 

 So there exists some \( n=k \) for which 

 \( \binom{1000}{n} \cdot \left(\frac{1}{5}\right)^n \)

 is as great as possible.

 particularly \( n \ge 165 \frac{5}{6} \)

 $\textbf{Incorporation of the Conjecture}$

 Every term in the binomial expansion of \( (1+0.2)^{1000} \) must be in the form 

 
 \[A_k = \binom{1000}{k}(0.2)^k\]

 where \( k \in \mathbb{Z}^+ : 0 \le k \le 1000 \).

 The maximal value occurs at \( k \ge 165 \frac{5}{6} \). However, this isn't an integer. The next greatest integer will occur at \( k = 166 \). Thus the value of \( k \) such that \( A_k \) is maximized is \(\boxed{\textbf{166}}\) \(\square\)

 





\textbf{Solution 2}

Let $0<x_{}^{}<1$. Then we may write $A_{k}^{}={N\choose k}x^{k}=\frac{N!}{k!(N-k)!}x^{k}=\frac{(N-k+1)!}{k!}x^{k}$. Taking logarithms in both sides of this last equation and using the well-known fact $\log(a_{}^{}b)=\log a + \log b$ (valid if $a_{}^{},b_{}^{}>0$), we have

 $\log(A_{k})=\log\left[\frac{(N-k+1)!}{k!}x^{k}\right]=\log\left[\prod_{j=1}^{k}\frac{(N-j+1)x}{j}\right]=\sum_{j=1}^{k}\log\left[\frac{(N-j+1)x}{j}\right]\, .$

 Now, $\log(A_{k}^{})$ keeps increasing with $k_{}^{}$ as long as the arguments $\frac{(N-j+1)x}{j}>1$ in each of the $\log\big[\;\big]$ terms (recall that $\log y_{}^{} <0$ if $0<y_{}^{}<1$). Therefore, the integer $k_{}^{}$ that we are looking for must satisfy $k=\Big\lfloor\frac{(N+1)x}{1+x}\Big\rfloor$, where $\lfloor z_{}^{}\rfloor$ denotes the greatest integer less than or equal to $z_{}^{}$. 

 In summary, substituting $N_{}^{}=1000$ and $x_{}^{}=0.2$ we finally find that $k=\boxed{166}$.

 





\textbf{Solution 3}

We know that once we have found the largest value of $k$, all values after $A_k$ are less than $A_k$. Therefore, we are looking for the smallest possible value such that:

 ${\frac{1}{5}}^k\cdot {{1000} \choose {k}}>{\frac{1}{5}}^{k+1}\cdot {{1000} \choose {k+1}}$

 Dividing by ${\frac{1}{5}}^k$ gives:

 ${1000\choose k}>{\frac{1}{5}}\cdot{1000\choose {k+1}}$.

 We can express these binomial coefficients as factorials.

 $\frac{1000!}{(1000-k)!\cdot(k)!}>{\frac{1}{5}}\cdot\frac{1000!}{(1000-k-1)!\cdot{(k+1)!}}$

 We note that the $1000!$ can cancel. Also, $(1000-k)!=(1000-k)(1000-k-1)!$. Similarly, $(k+1)!=(k+1)k!$.

 Canceling these terms yields,

 $\frac{1}{(1000-k)}>{\frac{1}{5}}\cdot\frac{1}{(k+1)}$

 Cross multiplying gives:

 $5k+5>1000-k \Rightarrow k>165.8$

 Therefore, since this identity holds for all values of $k>165.8$, the largest possible value of $k$ is $\boxed{166}$.

 





\textbf{Solution 4}

We know that $A_k$ will increase as $k$ increases until certain $k=m$, where $A_0 < A_1 < A_2 < \dots < A_{m-2} < A_{m-1} < A_m$ and  

 $A_m > A_{m+1} > A_{m+2} > \dots > A_{1000}.$

 Next, to change $A_{k-1}$ to $A_k$, we multiply $A_{k-1}$ by $\frac{1000-k+1}{5k}$. It follows that the numerator must be greater than the

 denominator or

 
 \[1000-k+1>5k\] 
 \[k<166.8\].

 We conclude that the answer is $\boxed{166}$.

 





\textbf{Solution 5}

Notice that the expansion is the largest the moment BEFORE the $nC_p < 5$ (this reasoning can probably be found in the other solutions; basically, if we have a number k and then k+1, the value is the largest when k+1 is larger than k, or in other words $nC_p*\frac{1}{5} > 1$)

 Say we have ${1000 \choose 5}$... this equals $\frac{1000*999*998*997*996}{5*4*3*2*1}$, and if we have k+1, then it's just going to be equivalent to multiplying this fraction by $\frac{995}{6}$. Notice that this fraction's numerator plus denominator is equal to $1001$. Calling the numerator x and the denominator y, we get that
 \[x+y = 1001\] 
 \[\frac{x}{y} > 5\] 
 \[x>5y\] 
 \[6y<1001\] 
 \[y<166.83333\] and since the denominator is what we are specifically choosing, or in other words what k is, we conclude k = $\boxed{166}$

 Note: It may be more intuitive to equate $x=5y$ and then use test points to figure out which side of the inequality to lie on.

 





\textbf{Solution 6}

Notice the relation between $A_m$ and $A_{m+1}$. We have that: $A_{m+1}  = A_m \cdot \frac{1}{5} \cdot \frac{1000-m}{m+1}$. This is true because from $A_m$ to $A_{m+1}$ we have to multiply by $\frac{1}{5}=0.2$ once,and then we must resolve the factorial issue. To do this, we must realize that 

 
 \[{1000 \choose m+1} = \frac{1000!}{(m+1)!(1000-m-1)!} = \frac{1000 \cdots (1001-m)(1000-m)}{(m+1)!}\]and that
 \[{1000 \choose m} = \frac{1000!}{m!(1000-m)!} = \frac{1000 \cdots (1001-m)}{m!}\]

 
 \[\Longrightarrow {1000 \choose m} \cdot \frac{1000-m}{m+1}\]

 So now, we must find out for which $m$ is $1000-m< 5 \cdot (m+1)$ (when this happens the numerator is less than the denominator for the fraction $\frac{1000-m}{m+1}$ so then we will have $A_m > A_{m+1}$) Then we find that for all $m > 165.833333$, the above inequality ($1000-m< 5 \cdot (m+1)$) holds, so then $m=165$ so $k=m+1= \boxed{166}$.

 




	\item (\textit{AIME I 1991}) How many real numbers $x^{}_{}$ satisfy the equation $\frac{1}{5}\log_2 x = \sin (5\pi x)$?


 
 \textbf{Answer: }159



\textbf{Solution}


                    \begin{figure}[H]
                        \centering
                        \includegraphics[width=0.5\textwidth]{Solution/AIME_1991_I_4_Solution_Figure_1.png}
                    \end{figure}
                    The \href{https://artofproblemsolving.com/wiki/index.php?title=Range}{range} of the \href{https://artofproblemsolving.com/wiki/index.php?title=Sine}{sine} function is $-1 \le y \le 1$. It is \href{https://artofproblemsolving.com/wiki/index.php?title=Periodic\_function}{periodic} (in this problem) with a period of $\frac{2}{5}$. 

 Thus, $-1 \le \frac{1}{5} \log_2 x \le 1$, and $-5 \le \log_2 x \le 5$. The solutions for $x$ occur in the domain of $\frac{1}{32} \le x \le 32$. When $x > 1$ the \href{https://artofproblemsolving.com/wiki/index.php?title=Logarithm}{logarithm} function returns a \href{https://artofproblemsolving.com/wiki/index.php?title=Positive}{positive} value; up to $x = 32$ it will pass through the sine curve. There are exactly 10 intersections of five periods (every two integral values of $x$) of the sine curve and another curve that is $< 1$, so there are $\frac{32}{2} \cdot 10 - 6 = 160 - 6 = 154$ values (the subtraction of 6 since all the “intersections” when $x < 1$ must be disregarded). When $y = 0$, there is exactly $1$ touching point between the two functions: $\left(\frac{1}{5},0\right)$. When $y < 0$ or $x < 1$, we can count $4$ more solutions. The solution is $154 + 1 + 4 = \boxed{159}$.

 





\textbf{Solution 2}

Notice that the equation is satisfied twice for every sine period (which is $\frac{2}{5}$), except in the sole case when the two equations equate to $0$. In that case, the equation is satisfied twice but only at the one instance when $y=0$. Hence, it is double-counted in our final solution, so we have to subtract it out. We then compute:$32 \cdot \frac{5}{2} \cdot 2 - 1 = \boxed {159}$

 




	\item (\textit{AIME I 1991}) Given a rational number, write it as a fraction in lowest terms and calculate the product of the resulting numerator and denominator. For how many rational numbers between 0 and 1 will  $20_{}^{}!$ be the resulting product? 


 
 \textbf{Answer: }128



\textbf{Solution}

If the fraction is in the form $\frac{a}{b}$, then $a < b$ and $\gcd(a,b) = 1$. There are 8 \href{https://artofproblemsolving.com/wiki/index.php?title=Prime\_number}{prime numbers} less than 20 ($2, 3, 5, 7, 11, 13, 17, 19$), and each can only be a factor of one of $a$ or $b$. There are $2^8$ ways of selecting some \href{https://artofproblemsolving.com/wiki/index.php?title=Combination}{combination} of numbers for $a$; however, since $a<b$, only half of them will be between $0 < \frac{a}{b} < 1$. Therefore, the solution is $\frac{2^8}{2} = \boxed{128}$.

 




	\item (\textit{AIME I 1991}) Suppose $r^{}_{}$ is a real number for which


 $\left\lfloor r + \frac{19}{100} \right\rfloor + \left\lfloor r + \frac{20}{100} \right\rfloor + \left\lfloor r + \frac{21}{100} \right\rfloor + \cdots + \left\lfloor r + \frac{91}{100} \right\rfloor = 546.$Find $\lfloor 100r \rfloor$. (For real $x^{}_{}$, $\lfloor x \rfloor$ is the greatest integer less than or equal to $x^{}_{}$.)


 
 \textbf{Answer: }743



\textbf{Solution 1}

There are $91 - 19 + 1 = 73$ numbers in the \href{https://artofproblemsolving.com/wiki/index.php?title=Sequence}{sequence}. Since the terms of the sequence can be at most $1$ apart, all of the numbers in the sequence can take one of two possible values. Since $\frac{546}{73} = 7 R 35$, the values of each of the terms of the sequence must be either $7$ or $8$. As the remainder is $35$, $8$ must take on $35$ of the values, with $7$ being the value of the remaining $73 - 35 = 38$ numbers. The 39th number is $\lfloor r+\frac{19 + 39 - 1}{100}\rfloor= \lfloor r+\frac{57}{100}\rfloor$, which is also the first term of this sequence with a value of $8$, so $8 \le r + \frac{57}{100} < 8.01$. Solving shows that $\frac{743}{100} \le r < \frac{744}{100}$, so $743\le 100r < 744$, and $\lfloor 100r \rfloor = \boxed{743}$.

 





\textbf{Solution 2 (Faster)}

Recall by Hermite's Identity that $\lfloor x\rfloor +\lfloor x+\frac{1}{n}\rfloor +...+\lfloor x+\frac{n-1}{n}\rfloor = \lfloor nx\rfloor$ for positive integers $n$, and real $x$. Similar to above, we quickly observe that the last 35 take the value of 8, and remaining first ones take a value of 7. So, $\lfloor r\rfloor \le ...\le \lfloor r+\frac{18}{100}\rfloor \le 7$ and $\lfloor r+1\rfloor \ge ...\ge \lfloor r+\frac{92}{100}\rfloor \ge 8$. We can see that $\lfloor r\rfloor +1=\lfloor r+1\rfloor$. Because $\lfloor r\rfloor$ is at most 7, and $\lfloor r+1\rfloor$ is at least 8, we can clearly see their values are $7$ and $8$ respectively. So,  $\lfloor r\rfloor = ... = \lfloor r+\frac{18}{100}\rfloor = 7$, and $\lfloor r+\frac{92}{100}\rfloor = ...= \lfloor r+1\rfloor = 8$. Since there are 19 terms in the former equation and 8 terms in the latter, our answer is $\lfloor nx\rfloor = 546+19\cdot 7+8\cdot 8=\boxed{743}$

 \textbf{Note} In the contest, you would just observe this mentally, and then calculate $546+19\cdot 7+8\cdot 8= 743$, hence the speed at which one can carry out this solution.

 





\textbf{Solution 3 (Algebra)}

Notice how all terms in this series is in the form \( \lfloor r + \frac{k}{100} \rfloor \), where \( 19 \le k \le 91 \). This then implies that there are a total of 73 terms in this series.

 If every term is the same, i.e., \( 73 \lfloor r + \frac{19}{100} \rfloor = 546 \), then \( \lfloor r + \frac{19}{100} \rfloor = \frac{546}{7} \approx 7 \). However, because this isn't exactly 7, this implies that the values in the series are mixed between 7 and 8, so we need to do a bit more work to find out the value of \( k \) such that \( \lfloor r + \frac{k}{100} \rfloor = 7 \) and \( \lfloor r + \frac{k+1}{100} \rfloor = 8 \).

 We see that if we have a values that return 7 and b values that return 8, then \( a+b=73 \).

 Additionally, because the total sum of all values is 546, we have \( 7a + 8b = 546 \).

 Solving the system returns \( b=35 \), so 35 terms return 8 and the others return 7.

 Then, it is notable that the last 35 terms of the series will return 8 (as they are the ones that maximize the value of r), and of course, we have \( k+1=91-35+1 =57 \). So \( k = 56 \).

 Therefore, the point of change occurs in the step from \( \lfloor r + \frac{56}{100} \rfloor \) to \( \lfloor r + \frac{57}{100} \rfloor \).

 We now use the definition of \( \lfloor r + \frac{57}{100} \rfloor = 8 \), because if \( \lfloor r + \frac{56}{100} \rfloor = 7 \), that would imply that \( \{r\} = \frac{100-57}{100} = \frac{43}{100} \) (as \( \lfloor \frac{56+43}{100} \rfloor = 0 \) and \( \lfloor \frac{57+43}{100} \rfloor = 1 \))

 Then, through the definition of \( \{r\} \), we see that \( r - \lfloor r \rfloor = \{r\} \), and because \( \lfloor r \rfloor = 7 \) and \( \{r\} = \frac{43}{100} \), we get \( r = \frac{743}{100} = 7.43 \).

 \( \lfloor 100r \rfloor = \lfloor 100 \cdot 7.43 \rfloor = \lfloor 743 \rfloor = \) $\boxed{743}$

 





\textbf{Solution 4 (Very Intuitive)}

Let $k = \lfloor r \rfloor \text{, for some integer } k \in \mathbb{Z}$.Plugging in into the expression simplifies into,
 \[73k = 546 \Rightarrow k = \frac{546}{73} = 7.479...\]Since $k$ must be an integer, this can't be a solution. Instead, in order for it to be a solution, the numerator must be divisible by $73$. Notice that the closest multiple of $73$ less than $546$ is $511$, and $546 - 511 = 35$. Hence, $k$ will have an integer solution if we can achieve the expression $73k + 35 = 546$.

 The last term of the sum is $\lfloor r + \frac{91}{100} \rfloor$ , and in order for that floor function to equal $k + 1$, $r$ must be in the form 
 \[r = k + 0.09 + 0.01x \text{, for some real number } x \in \mathbb{R}, 0 \le x < 1\]This needs to go on until the $35$th last term $\left(\text{from } \left\lfloor r + \frac{57}{100} \right\rfloor \text{ to } \left\lfloor r + \frac{91}{100} \right\rfloor\right)$ in order to achieve $73k + 35 = 546$. For this to happen, the last 35 terms must each equal $k + 1$, hence $r$ must satisfy $r = k + 0.43 + 0.01x$.Hence, we achieve,
 \[73k + 35 = 546\]
 \[73k = 511 \Rightarrow k = 7\]Which means,
 \[r = k + 0.43 + 0.01x \Rightarrow r = 7.43 + 0.01x\]Hence, $\lfloor 100r \rfloor = \lfloor 743 + x \rfloor = \boxed{743}$

 




	\item (\textit{AIME I 1991}) Find $A^2_{}$, where $A^{}_{}$ is the sum of the absolute values of all roots of the following equation:


 $x = \sqrt{19} + \frac{91}{{\sqrt{19}+\frac{91}{{\sqrt{19}+\frac{91}{{\sqrt{19}+\frac{91}{{\sqrt{19}+\frac{91}{x}}}}}}}}}$
 \textbf{Answer: }383



\textbf{Solution 1}

$x=\sqrt{19}+\underbrace{\frac{91}{\sqrt{19}+\frac{91}{\sqrt{19}+\frac{91}{\sqrt{19}+\frac{91}{\sqrt{19}+\frac{91}{x}}}}}}_{x}$

 $x=\sqrt{19}+\frac{91}{x}$

 $x^2=x\sqrt{19}+91$

 $x^2-x\sqrt{19}-91 = 0$

 $\left.\begin{array}{l}x_1=\frac{\sqrt{19}+\sqrt{383}}{2}\\\\x_2=\frac{\sqrt{19}-\sqrt{383}}{2}\end{array}\right\}A=|x_1|+|x_2|\Rightarrow\sqrt{383}$

 $A^2=\boxed{383}$

 \textbf{Proof of the first step} Think about the set for $x$. The first element of this set is $\sqrt{19}+\frac{91}{x}$. The next is $\sqrt{19}+\frac{91}{\sqrt{19}+\frac{91}{x}}$, and so on so forth. For simplicity, $\mathcal{A} = \{\sqrt{19}+\frac{91}{x},\sqrt{19}+\frac{91}{\sqrt{19}+\frac{91}{x}}, \dots \}$. Then, notice that the set $\mathcal{B}$ containing first term $\sqrt{19}+\frac{91}{\sqrt{19}+\frac{91}{x}}$ is simply substituting the definition of $x$ into the first term of $\mathcal{A}$. Therefore, we can simply multiply each element of $\mathcal{B}$ to obtain the same element of $\mathcal{A}$. Then, we can create a set $\mathcal{C}$ such that $x=\sqrt{19}+\frac{91}{\sqrt{19}+\frac{91}{\sqrt{19}+\frac{91}{\sqrt{19}+\frac{91}{\sqrt{19}+\frac{91}{x}}}}}$ is equal to $x=\sqrt{19}+\frac{91}{x}$.

 The motivation for this arises through showing a bijection between two infinite sets, and that they have the same number of elements.

 Hope this helped.

 





\textbf{Solution 2}

Let $f(x) = \sqrt{19} + \frac{91}{x}$. Then $x = f(f(f(f(f(x)))))$, from which we realize that $f(x) = x$. This is because if we expand the entire expression, we will get a fraction of the form $\frac{ax + b}{cx + d}$ on the right hand side, which makes the equation simplify to a quadratic. As this quadratic will have two roots, they must be the same roots as the quadratic $f(x)=x$.

 The given finite expansion can then be easily seen to reduce to the \href{https://artofproblemsolving.com/wiki/index.php?title=Quadratic\_equation}{quadratic equation}, $x_{}^{2}-\sqrt{19}x-91=0$. The solutions are $x_{\pm}^{}=$$\frac{\sqrt{19}\pm\sqrt{383}}{2}$. Therefore, $A_{}^{}=\vert x_{+}\vert+\vert x_{-}\vert=\sqrt{383}$. We conclude that $A_{}^{2}=\boxed{383}$.

 




	\item (\textit{AIME I 1991}) For how many real numbers $a^{}_{}$ does the quadratic equation $x^2 + ax^{}_{} + 6a=0$ have only integer roots for $x^{}_{}$?


 
 \textbf{Answer: }010



\textbf{Solution 1}

Let $x^2 + ax + 6a = (x - s)(x - r)$. Vieta's yields $s + r = - a, sr = 6a$.
 \begin{eqnarray*}sr + 6s + 6r &=& 0\\ sr + 6s + 6r + 36 &=& 36\\ (s + 6)(r + 6) &=& 36 \end{eqnarray*}

 \href{https://artofproblemsolving.com/wiki/index.php?title=Without\_loss\_of\_generality}{Without loss of generality} let $r \le s$.

 The possible values of $(r + 6,s + 6)$ are:$( - 36, - 1),( - 18, - 2),( - 12, - 3),( - 9, - 4),( - 6, - 6),(1,36),(2,18),(3,12),(4,9),(6,6)$$\Rightarrow \boxed{10}\ \text{values of } a$.

 





\textbf{Solution 2}

By \href{https://artofproblemsolving.com/wiki/index.php?title=Vieta\%27s\_formulas}{Vieta's formulas}, $x_1 + x_2 = -a$ where $x_1, x_2$ are the roots of the quadratic, and since $x_1,x_2$ are integers, $a$ must be an integer. Applying the \href{https://artofproblemsolving.com/wiki/index.php?title=Quadratic\_formula}{quadratic formula}, 

 
 \[x = \frac{-a \pm \sqrt{a^2 - 24a}}{2}\]

 Since $-a$ is an integer, we need $\sqrt{a^2-24a}$ to be an integer (let this be $b$): $b^2 = a^2 - 24a$. \href{https://artofproblemsolving.com/wiki/index.php?title=Completing\_the\_square}{Completing the square}, we get

 
 \[(a - 12)^2 = b^2 + 144\]

 Which implies that $b^2 + 144$ is a \href{https://artofproblemsolving.com/wiki/index.php?title=Perfect\_square}{perfect square} also (let this be $c^2$). Then 

 
 \[c^2 - b^2 = 144 \Longrightarrow (c+b)(c-b) = 144\]

 The pairs of factors of $144$ are $(\pm1,\pm144),( \pm 2, \pm 72),( \pm 3, \pm 48),( \pm 4, \pm 36),( \pm 6, \pm 24),( \pm 8, \pm 18),( \pm 9, \pm 16),( \pm 12, \pm 12)$; since $c$ is the \href{https://artofproblemsolving.com/wiki/index.php?title=Mean}{average} of each respective pair and is also an integer, the pairs that work must have the same \href{https://artofproblemsolving.com/wiki/index.php?title=Parity}{parity}. Thus we get $\boxed{10}$ pairs (counting positive and negative) of factors that work, and substituting them backwards show that they all work.

 




	\item (\textit{AIME I 1991}) Suppose that $\sec x+\tan x=\frac{22}7$ and that $\csc x+\cot x=\frac mn,$ where $\frac mn$ is in lowest terms.  Find $m+n^{}_{}.$


 
 \textbf{Answer: }044



\textbf{Solution 1}

Use the two \href{https://artofproblemsolving.com/wiki/index.php?title=Trigonometric\_identities\#Pythagorean\_Identities}{trigonometric Pythagorean identities} $1 + \tan^2 x = \sec^2 x$ and $1 + \cot^2 x = \csc^2 x$. 

 If we square the given $\sec x = \frac{22}{7} - \tan x$, we find that 

 
 \begin{align*} \sec^2 x &= \left(\frac{22}7\right)^2 - 2\left(\frac{22}7\right)\tan x + \tan^2 x \\ 1 &= \left(\frac{22}7\right)^2 - \frac{44}7 \tan x \end{align*}


 This yields $\tan x = \frac{435}{308}$. 

 Let $y = \frac mn$. Then squaring, 

 
 \[\csc^2 x = (y - \cot x)^2 \Longrightarrow 1 = y^2 - 2y\cot x.\] 

 Substituting $\cot x = \frac{1}{\tan x} = \frac{308}{435}$ yields a \href{https://artofproblemsolving.com/wiki/index.php?title=Quadratic\_equation}{quadratic equation}: $0 = 435y^2 - 616y - 435 = (15y - 29)(29y + 15)$. It turns out that only the \href{https://artofproblemsolving.com/wiki/index.php?title=Positive}{positive} root will work, so the value of $y = \frac{29}{15}$ and $m + n = \boxed{044}$.

 Note: The problem is much easier computed if we consider what $\sec (x)$ is, then find the relationship between $\sin( x)$ and $cos (x)$ (using $\tan (x) = \frac{435}{308}$, and then computing $\csc x + \cot x$ using $1/\sin x$ and then the reciprocal of $\tan x$.

 





\textbf{Solution 2}

Recall that $\sec^2 x - \tan^2 x = 1$, from which we find that $\sec x - \tan x = 7/22$. Adding the equations

 
 \begin{eqnarray*} \sec x + \tan x & = & 22/7 \\ \sec x - \tan x & = & 7/22\end{eqnarray*}

 together and dividing by 2 gives $\sec x = 533/308$, and subtracting the equations and dividing by 2 gives $\tan x = 435/308$. Hence, $\cos x = 308/533$ and $\sin x = \tan x \cos x = (435/308)(308/533) = 435/533$. Thus, $\csc x = 533/435$ and $\cot x = 308/435$. Finally,

 
 \[\csc x + \cot x = \frac {841}{435} = \frac {29}{15},\]

 so $m + n = 044$.

 





\textbf{Solution 3 (least computation)}

By the given,$\frac {1}{\cos x} + \frac {\sin x}{\cos x} = \frac {22}{7}$ and$\frac {1}{\sin x} + \frac {\cos x}{\sin x} = k$.

 Multiplying the two, we have

 
 \[\frac {1}{\sin x \cos x} + \frac {1}{\sin x} + \frac {1}{\cos x} + 1 = \frac {22}{7}k\]

 Subtracting both of the two given equations from this, and simpliyfing with the identity $\frac {\sin x}{\cos x} + \frac {\cos x}{\sin x} = \frac{\sin ^2 x + \cos ^2 x}{\sin x \cos x} = \frac {1}{\sin x \cos x}$, we get 

 
 \[1 = \frac {22}{7}k - \frac {22}{7} - k.\]

 Solving yields $k = \frac {29}{15}$, and $m+n = 044$

 





\textbf{Solution 4}

Make the substitution $u = \tan \frac x2$ (a substitution commonly used in calculus). By the half-angle identity for tangent, $\tan \frac x2 = \frac{\sin x}{1+\cos x}$, so $\csc x + \cot x = \frac{1+\cos x}{\sin x} = \frac1u = \frac mn$. Also, we have $\sec x + \tan x = \frac{1 + \sin x}{\cos x}.$ Now note the following:

 
 \begin{align*}\sin x &= \frac{2u}{1+u^2}\\ \cos x &= \frac{1-u^2}{1+u^2}\end{align*}


 Plugging these into our equality gives:

 
 \[\frac{1+\frac{2u}{1+u^2}}{\frac{1-u^2}{1+u^2}} = \frac{22}7\]

 This simplifies to $\frac{1+u}{1-u} = \frac{22}7$, and solving for $u$ gives $u = \frac{15}{29}$, and $\frac mn = \frac{29}{15}$. Finally, $m+n = 044$.

 





\textbf{Solution 5}

We are given that $\frac{1+\sin x}{\cos x}=\frac{22}7\implies\frac{1+\sin x}{\cos x}\cdot\frac{1-\sin x}{1-\sin x}=\frac{1-\sin^2x}{\cos x(1-\sin x)}=\frac{\cos^2x}{\cos x(1-\sin x)}$$=\frac{\cos x}{1-\sin x}$, or equivalently, $\cos x=\frac{7+7\sin x}{22}=\frac{22-22\sin x}7\implies\sin x=\frac{22^2-7^2}{22^2+7^2}$$\implies\cos x=\frac{2\cdot22\cdot7}{22^2+7^2}$. Note that what we want is just $\frac{1+\cos x}{\sin x}=\frac{1+\frac{2\cdot22\cdot7}{22^2+7^2}}{\frac{22^2-7^2}{22^2+7^2}}=\frac{22^2+7^2+2\cdot22\cdot7}{22^2-7^2}=\frac{(22+7)^2}{(22-7)(22+7)}=\frac{22+7}{22-7}$$=\frac{29}{15}\implies m+n=29+15=\boxed{044}$.

 





\textbf{Solution 6}

Assign a right triangle with angle $x$, hypotenuse $c$, adjacent side $a$, and opposite side $b$.Then, through the given information above, we have that..

 $\frac{c}{a}+\frac{b}{a}=\frac{22}{7}\implies \frac{c+b}{a}=\frac{22}{7}$

 $\frac{c}{b}+\frac{a}{b}=\frac{m}{n}\implies \frac{a+c}{b}=\frac{m}{n}$

 Hence, because similar right triangles can be scaled up by a factor, we can assume that this particular right triangle is indeed in simplest terms.

 Hence, $a=7$, $b+c=22$

 Furthermore, by the Pythagorean Theorem, we have that

 $a^2+b^2=c^2\implies 49+b^2=c^2$

 Solving for $c$ in the first equation and plugging in into the second equation...

 $49+b^2=(22-b)^2\implies 49+b^2=484-44b+b^2\implies 44b=435\implies b=\frac{435}{44}$

 Hence, $c=22-\frac{435}{44}=\frac{533}{44}$

 Now, we want $\frac{a+c}{b}$

 Plugging in, we find the answer is $\frac{\frac{7\cdot{44}}{44}+\frac{533}{44}}{\frac{435}{44}}=\frac{841}{435}=\frac{29}{15}$

 Hence, the answer is $29+15=\boxed{044}$

 





\textbf{Solution 7}

We know that $\sec(x) = \frac{h}{a}$ and that $\tan(x) = \frac{o}{a}$ where $h$, $a$, $o$ represent the hypotenuse, adjacent, and opposite (respectively) to angle $x$ in a right triangle. Thus we have that $\sec(x) + \tan(x) = \frac{h+o}{a}$. We also have that $\csc(x) + \cot(x) = \frac{h}{o} + \frac{a}{o} = \frac{h+a}{o}$. Set $\sec(x) + \tan(x) = \alpha$ and csc(x)+cot(x) = $\beta$. Then, notice that $\alpha + \beta = \frac{h+o}{a} + \frac{h+a}{o} = \frac{oh+ah+o^2 + a^2}{oa} = \frac{h(o+a+h)}{oa}$ ( This is because of the Pythagorean Theorem, recall $o^2 +a^2 = h^2$). But then notice that $\alpha \cdot \beta = \frac{(o+h)(a+h)}{oa} = \frac{oa +oh +ha +h^2}{oa} = 1+ \frac{h(o+a+h)}{oa} = 1+ \alpha + \beta$. From the information provided in the question, we can substitute $\alpha$ for $\frac{22}{7}$. Thus, $\frac{22 \beta}{7}= \beta + \frac{29}{7} \Longrightarrow 22 \beta = 7 \beta + 29 \Longrightarrow 15 \beta = 29 \Longrightarrow \beta = \frac{29}{15}$. Since, essentially we are asked to find the sum of the numerator and denominator of $\beta$, we have $29 + 15 = \boxed{044}$.

 


 





\textbf{Solution 8}

Firstly, we write $\sec x+\tan x=a/b$ where $a=22$ and $b=7$. This will allow us to spot factorable expressions later. Now, since $\sec^2x-\tan^2x=1$, this gives us 
 \[\sec x-\tan x=\frac{b}{a}\] Adding this to our original expressions gives us 
 \[2\sec x=\frac{a^2+b^2}{ab}\] or 
 \[\cos x=\frac{2ab}{a^2+b^2}\] Now since $\sin^2x+\cos^2x=1$, $\sin x=\sqrt{1-\cos^2x}$ So we can write 
 \[\sin x=\sqrt{1-\frac{4a^2b^2}{(a^2+b^2)^2}}\] Upon simplification, we get 
 \[\sin x=\frac{a^2-b^2}{a^2+b^2}\] We are asked to find $1/\sin x+\cos x/\sin x$ so we can write that as 
 \[\csc x+\cot x=\frac{1}{\sin x}+\frac{\cos x}{\sin x}\] 
 \[\csc x+\cot x=\frac{a^2+b^2}{a^2-b^2}+\frac{2ab}{a^2+b^2}\frac{a^2+b^2}{a^2-b^2}\] 
 \[\csc x+\cot x=\frac{a^2+b^2+2ab}{a^2-b^2}\] 
 \[\csc x+\cot x=\frac{(a+b)^2}{(a-b)(a+b)}\] 
 \[\csc x+\cot x=\frac{a+b}{a-b}\] Now using the fact that $a=22$ and $b=7$ yields, 
 \[\csc x+\cot x=\frac{29}{15}=\frac{p}{q}\] so $p+q=15+29=\boxed{44}$

 





\textbf{Solution 9}

Rewriting $\sec{x}$ and $\tan{x}$ in terms of $\sin{x}$ and $\cos{x}$, we know that $\frac{1+\sin{x}}{\cos{x}}=\frac{22}{7}.$

 Clearing fractions, 
 \[22\cos{x}=7+7\sin{x}.\]

 Squaring to get an expression in terms of $\sin^2{x}$ and $\cos^2{x}$, 
 \[484\cos^2{x}=49+49\sin^2{x}+98\sin{x}.\]

 Substituting $\cos^2{x}=1-\sin^2{x},$

 
 \[484(1-\sin^2{x})=49+49\sin^2{x}+98\sin{x}.\]

 Expanding then collecting terms yields a quadratic in $\sin{x}:$

 
 \[533\sin^2{x}+98\sin{x}-435=0.\]

 To make calculations easier, let $y=\sin{x}.$ 

 
 \[533y^2+98y-435=0.\]

 Upon inspection, $y=-1$ is a root. Dividing by $y+1$, 

 
 \[533y^2+98y-435=(533y-435)(y+1).\]

 Substituting $y=\sin{x},$ we see that $\sin{x}=-1$ doesn't work, as $\cos{x}=0$, leaving $\tan{x}$ undefined. 

 We conclude that $\sin{x}=\frac{435}{533}.$

 Since $\sin^2{x}+\cos^2{x}=1,$

 
 \[\cos{x}=\pm \sqrt{\frac{533^2-435^2}{533^2}}.\]
 \[=\pm \frac{308}{533}.\]

 After checking via the given equation, we know that only the positive solution works. 

 Therefore,

 
 \[\csc{x}+\cot{x}=\frac{1}{\sin{x}}+\frac{\cos{x}}{\sin{x}}\]
 \[=\frac{533}{435}+\frac{308}{435}\]
 \[=\frac{29}{15}=\frac{m}{n}.\]

 Adding $m$ and $n$, our answer is $\boxed{\textbf{044}}.$

 -Benedict T (countmath1)

 





\textbf{Solution 10}

First note that
 \[\csc{x} + \cot{x} = \frac{1+\cos{x}}{\sin{x}}=\frac{2\cos^2{\frac{x}{2}}}{2\sin{\frac{x}{2}}\cos{\frac{x}{2}}}=\cot{\frac{x}{2}},\]as well as
 \[\frac{22}{7}=\sec{x} + \tan{x} = \frac{1+\sin{x}}{\cos{x}} = \frac{(\sin{\frac{x}{2}} + \cos{\frac{x}{2}})^2}{\cos^2{\frac{x}{2}} - \sin^2{\frac{x}{2}}}=\frac{\sin{\frac{x}{2}} + \cos{\frac{x}{2}}}{\cos{\frac{x}{2}} -\sin{\frac{x}{2}}} = \frac{\cot{\frac{x}{2}}+1}{\cot{\frac{x}{2}}-1}.\]Solving this, we have $\cot{x}=\frac{29}{15}.$ Thus, the answer is $\boxed{044}.$

 




	\item (\textit{AIME I 1991}) Two three-letter strings, $aaa^{}_{}$ and $bbb^{}_{}$, are transmitted electronically. Each string is sent letter by letter. Due to faulty equipment, each of the six letters has a 1/3 chance of being received incorrectly, as an $a^{}_{}$ when it should have been a $b^{}_{}$, or as a $b^{}_{}$ when it should be an $a^{}_{}$. However, whether a given letter is received correctly or incorrectly is independent of the reception of any other letter. Let $S_a^{}$ be the three-letter string received when $aaa^{}_{}$ is transmitted and let $S_b^{}$ be the three-letter string received when $bbb^{}_{}$ is transmitted. Let $p$ be the probability that $S_a^{}$ comes before $S_b^{}$ in alphabetical order. When $p$ is written as a fraction in lowest terms, what is its numerator?


 
 \textbf{Answer: }532



\textbf{Solution 1}

Let us make a chart of values in alphabetical order, where $P_a,\ P_b$ are the probabilities that each string comes from $aaa$ and $bbb$ multiplied by $27$, and $S_b$ denotes the partial sums of $P_b$ (in other words, $S_b = \sum_{n=1}^{b} P_b$):
 \[\begin{array}{|r||r|r|r|} \hline \text{String}&P_a&P_b&S_b\\ \hline aaa & 8 & 1 & 1 \\ aab & 4 & 2 & 3 \\ aba & 4 & 2 & 5 \\ abb & 2 & 4 & 9 \\ baa & 4 & 2 & 11 \\ bab & 2 & 4 & 15 \\ bba & 2 & 4 & 19 \\ bbb & 1 & 8 & 27 \\ \hline \end{array}\]

 The probability is $p=\sum P_a \cdot (27 - S_b)$, so the answer turns out to be $\frac{8\cdot 26 + 4 \cdot 24 \ldots 2 \cdot 8 + 1 \cdot 0}{27^2} = \frac{532}{729}$, and the solution is $\boxed{532}$.

 





\textbf{Solution 2}

Let $S(a,n)$ be the $n$th letter of string $S(a)$.Compare the first letter of the string $S(a)$ to the first letter of the string $S(b)$.There is a $(2/3)^2=4/9$ chance that $S(a,1)$ comes before $S(b,1)$.There is a $2(1/3)(2/3)=4/9$ that $S(a,1)$ is the same as $S(b,1)$.

 If $S(a,1)=S(b,1)$, then you do the same for the second letters of the strings.  But you have to multiply the $4/9$ chance that $S(a,2)$ comes before $S(b,2)$ as there is a $4/9$ chance we will get to this step.

 Similarly, if $S(a,2)=S(b,2)$, then there is a $(4/9)^3$ chance that we will get to comparing the third letters and that $S(a)$ comes before $S(b)$.

 So we have $p=(4/9)+(4/9)^2+(4/9)^3=4/9+16/81+64/729=532/729$. Therefore, the answer is $\boxed{532}$. 

 





\textbf{Solution 3}

Consider $n$ letter strings instead. If the first letters all get transmitted correctly, then the $a$ string will be first. Otherwise, the only way is for both of the first letters to be the same, and then we consider the next $n-1$ letter string following the first letter. This easily leads to a recursion: $p_n=\frac23\cdot\frac23+2\cdot\frac23\cdot\frac13p_{n-1}=\frac49+\frac49p_{n-1}$. Clearly, $p_0=0\implies p_1=\frac49\implies p_2=\frac{52}{81}\implies p_3=\frac{532}{729}$. Therefore, the answer is $\boxed{532}$.

 





\textbf{Solution 4 (a more explicit Solution 3)}

The probability that $S_a$ will take the form $a$ \_ \_ and that $S_b$ will take the form $b$ \_ \_ is $\frac{2}{3}\cdot\frac{2}{3} = \frac{4}{9}$. Then, the probability that both $S_a$ and $S_b$ will share the same first digit is $2\cdot\frac{2}{3}\cdot\frac{1}{3} = \frac{4}{9}$. Now if the first digits of either sequence are the same, then we must now consider these same probabilities for the second letter of each sequence. The probability that when the first two letters of both sequences are the same, that the second letter of $S_a$ is $a$ and that the second letter of $S_b$ is $b$ is $\frac{2}{3}\cdot\frac{2}{3}=\frac{4}{9}$. Similarly, the probability that when the first two letters of both sequences are the same, that the second set of letters in both sets of sequences are the same is $2\cdot\frac{2}{3}\cdot\frac{1}{3} = \frac{4}{9}$. Now, if the last case is true then the probability that $S_a$ precedes $S_b$ is $\frac{2}{3}\cdot\frac{2}{3} = \frac{4}{9}$. Therefore the total probability would be: $\frac{4}{9} + \frac{4}{9}\left(\frac{4}{9} + \frac{4}{9}\left(\frac{4}{9}\right)\right) = \frac{4}{9}+\frac{4}{9}\left(\frac{52}{81}\right) = \frac{4}{9} + \frac{208}{729} = \frac{532}{729}$. Therefore the answer is $\boxed{532}$.

 




	\item (\textit{AIME I 1991}) Twelve congruent disks are placed on a circle $C^{}_{}$ of radius 1 in such a way that the twelve disks cover $C^{}_{}$, no two of the disks overlap, and so that each of the twelve disks is tangent to its two neighbors. The resulting arrangement of disks is shown in the figure below.  The sum of the areas of the twelve disks can be written in the form $\pi(a-b\sqrt{c})$, where $a,b,c^{}_{}$ are positive integers and $c^{}_{}$ is not divisible by the square of any prime. Find $a+b+c^{}_{}$.


 
                    \begin{figure}[H]
                        \centering
                        \includegraphics[width=0.5\textwidth]{Problem/AIME_1991_I_11_Problem_Figure_1.png}
                    \end{figure}
                    


 
 \textbf{Answer: }135



\textbf{Solution}

We wish to find the radius of one circle, so that we can find the total area.

 Notice that for them to contain the entire circle, each pair of circles must be tangent on the larger circle. Now consider two adjacent smaller circles. This means that the line connecting the radii is a segment of length $2r$ that is tangent to the larger circle at the midpoint of the two centers. Thus, we have essentially have a regular dodecagon whose vertices are the centers of the smaller triangles circumscribed about a circle of radius $1$.

 We thus know that the \href{https://artofproblemsolving.com/wiki/index.php?title=Apothem}{apothem} of the \href{https://artofproblemsolving.com/wiki/index.php?title=Dodecagon}{dodecagon} is equal to $1$. To find the side length, we make a triangle consisting of a vertex, the midpoint of a side, and the center of the dodecagon, which we denote $A, M,$ and $O$ respectively. Notice that $OM=1$, and that $\triangle OMA$ is a right triangle with \href{https://artofproblemsolving.com/wiki/index.php?title=Hypotenuse}{hypotenuse} $OA$ and $m \angle MOA = 15^\circ$. Thus $AM = (1) \tan{15^\circ} = 2 - \sqrt {3}$, which is the radius of one of the circles. (If you don't already know the value of $\tan 15^{\circ}$, it's straightforward to calculate by writing it as $\tan\left(45^{\circ}-30^{\circ}\right)$ and then using the tangent addition/subtraction formula.) The area of one circle is thus $\pi(2 - \sqrt {3})^{2} = \pi (7 - 4 \sqrt {3})$, so the area of all $12$ circles is $\pi (84 - 48 \sqrt {3})$, giving an answer of $84 + 48 + 3 = \boxed{135}$.

 \textbf{Helpful Diagram} 
                    \begin{figure}[H]
                        \centering
                        \includegraphics[width=0.5\textwidth]{Solution/AIME_1991_I_11_Solution_Figure_1.png}
                    \end{figure}
                    

 \_Diagram by 1-1 is 3\_

 




	\item (\textit{AIME I 1991}) Rhombus $PQRS^{}_{}$ is inscribed in rectangle $ABCD^{}_{}$ so that vertices $P^{}_{}$, $Q^{}_{}$, $R^{}_{}$, and $S^{}_{}$ are interior points on sides $\overline{AB}$, $\overline{BC}$, $\overline{CD}$, and $\overline{DA}$, respectively. It is given that $PB^{}_{}=15$, $BQ^{}_{}=20$, $PR^{}_{}=30$, and $QS^{}_{}=40$. Let $m/n^{}_{}$, in lowest terms, denote the perimeter of $ABCD^{}_{}$. Find $m+n^{}_{}$.


 
 \textbf{Answer: }677



\textbf{Solution 1}

Let $O$ be the center of the rhombus. Via \href{https://artofproblemsolving.com/wiki/index.php?title=Parallel}{parallel} sides and alternate interior angles, we see that the opposite \href{https://artofproblemsolving.com/wiki/index.php?title=Triangle}{triangles} are \href{https://artofproblemsolving.com/wiki/index.php?title=Congruent}{congruent} ($\triangle BPQ \cong \triangle DRS$, $\triangle APS \cong \triangle CRQ$). Quickly we realize that $O$ is also the center of the rectangle. 

 By the \href{https://artofproblemsolving.com/wiki/index.php?title=Pythagorean\_Theorem}{Pythagorean Theorem}, we can solve for a side of the rhombus; $PQ = \sqrt{15^2 + 20^2} = 25$. Since the \href{https://artofproblemsolving.com/wiki/index.php?title=Diagonal}{diagonals} of a rhombus are \href{https://artofproblemsolving.com/wiki/index.php?title=Perpendicular\_bisector}{perpendicular bisectors}, we have that $OP = 15, OQ = 20$. Also, $\angle POQ = 90^{\circ}$, so quadrilateral $BPOQ$ is \href{https://artofproblemsolving.com/wiki/index.php?title=Cyclic\_quadrilateral}{cyclic}. By \href{https://artofproblemsolving.com/wiki/index.php?title=Ptolemy\%27s\_Theorem}{Ptolemy's Theorem}, $25 \cdot OB = 20 \cdot 15 + 15 \cdot 20 = 600$.

 By similar logic, we have $APOS$ is a cyclic quadrilateral. Let $AP = x$, $AS = y$. The Pythagorean Theorem gives us $x^2 + y^2 = 625\quad \mathrm{(1)}$. Ptolemy’s Theorem gives us $25 \cdot OA = 20x + 15y$. Since the diagonals of a rectangle are equal, $OA = \frac{1}{2}d = OB$, and $20x + 15y = 600\quad \mathrm{(2)}$. Solving for $y$, we get $y = 40 - \frac 43x$. Substituting into $\mathrm{(1)}$, 

 
 \begin{eqnarray*}x^2 + \left(40-\frac 43x\right)^2 &=& 625\\ 5x^2 - 192x + 1755 &=& 0\\ x = \frac{192 \pm \sqrt{192^2 - 4 \cdot 5 \cdot 1755}}{10} &=& 15, \frac{117}{5}\end{eqnarray*}

 We reject $15$ because then everything degenerates into squares, but the condition that $PR \neq QS$ gives us a \href{https://artofproblemsolving.com/wiki/index.php?title=Contradiction}{contradiction}. Thus $x = \frac{117}{5}$, and backwards solving gives $y = \frac{44}5$. The perimeter of $ABCD$ is $2\left(20 + 15 + \frac{117}{5} + \frac{44}5\right) = \frac{672}{5}$, and $m + n = \boxed{677}$.

 





\textbf{Solution 2}

From above, we have $OB = 24$ and $BD = 48$. Returning to $BPQO,$ note that $\angle PQO\cong \angle PBO \cong ABD.$ Hence, $\triangle ABD \sim \triangle OQP$ by $AA$ \href{https://artofproblemsolving.com/wiki/index.php?title=Similar\_triangle&action=edit&redlink=1}{similarity}. From here, it's clear that
 \[\frac {AD}{BD} = \frac {OP}{PQ}\implies \frac {AD}{48} = \frac {15}{25}\implies AD = \frac {144}{5}.\]Similarly,
 \[\frac {AB}{BD} = \frac {IQ}{PQ}\implies \frac {AB}{48} = \frac {20}{25}\implies AB = \frac {192}{5}.\]Therefore, the perimeter of rectangle $ABCD$ is $2(AB + AD) = 2\left(\frac {192}{5} + \frac {144}{5}\right) = \frac {672}{5}.$

 





\textbf{Solution 3}

The triangles $QOB,OBC$ are \href{https://artofproblemsolving.com/wiki/index.php?title=Isosceles\_triangle}{isosceles}, and \href{https://artofproblemsolving.com/wiki/index.php?title=Similar\_triangles}{similar} (because they have $\angle QOB = \angle OBC$).

 Hence $\frac {BQ}{OB} = \frac {OB}{BC} \Rightarrow OB^2 = BC \cdot BQ$.

 The length of $OB$ could be found easily from the area of $BPQ$:

 
 \[BP \cdot BQ = \frac {OB}{2} \cdot PQ \Rightarrow OB = \frac {2BP\cdot BQ}{PQ} \Rightarrow OB = 24\]
 \[OB^2 = BC \cdot BQ \Rightarrow 24^2 = (20 + CQ) \cdot 20 \Rightarrow CQ = \frac {44}{5}\]

 From the right triangle $CRQ$ we have $RC^2 = 25^2 - \left(\frac {44}{5}\right)^2\Rightarrow RC = \frac {117}{5}$. We could have also defined a similar formula: $OB^2 = BP \cdot BA$, and then we found $AP$, the segment $OB$ is tangent to the circles with diameters $AO,CO$.

 The perimeter is $2(PB + BQ + QC + CR) = 2\left(15 + 20 + \frac {44 + 117}{5}\right) = \frac {672}{5}\Rightarrow m+n=677$.

 





\textbf{Solution 4}

For convenience, let $\angle PQS = \theta$. Since the opposite triangles are congruent we have that $\angle BQR = 3\theta$, and therefore $\angle QRC = 3\theta - 90$. Let $QC = a$, then we have $\sin{(3\theta - 90)} = \frac {a}{25}$, or $- \cos{3\theta} = \frac {a}{25}$. Expanding with the formula $\cos{3\theta} = 4\cos^3{\theta} - 3\cos{\theta}$, and since we have $\cos{\theta} = \frac {4}{5}$, we can solve for $a$. The rest then follows similarily from above.

 





\textbf{Solution 5}

We can just find coordinates of the points.  After drawing a picture, we can see 4 congruent right triangles with sides of $15,\ 20,\ 25$, namely triangles $DSR, OSR, OQP,$ and $BQP$.

 Let the points of triangle $DSR$ be $(0,0)\ (0,20)\ (15,0)$. Let point $E$ be on $\overline{SR}$, such that $SE = 16$ and $ER = 9$. Triangle $DSR$ can be split into two similar 3-4-5 right triangles, $ESD$ and $EDR$. By the Pythagorean Theorem, point $D$ is $12$ away from point $E$. Repeating the process, if we break down triangle $DER$ into two more similar triangles, we find that point $E$ is at $(9.6, 7.2)$.

 By reflecting point $D = (0,0)$ over point $E = (9.6, 7.2)$, we get point $O = (19.2, 14.4)$. By reflecting point $D$ over point $O$, we get point $B = (38.4, 28.8)$. Thus, the perimeter is equal to $(38.4 + 28.8)\times 2 = \frac {672}{5}$, making the final answer $672+5 = 677$.

 





\textbf{Solution 6}

We can just use areas. Let $AP = b$ and $AS = a$. $a^2 + b^2 = 625$. Also, we can add up the areas of all 8 right triangles and let that equal the total area of the rectangle, $(a+20)(b+15)$. This gives $3a + 4b = 120$. Solving this system of equation gives $\frac{44}{5} = a$, $\frac{117}{5} = b$, from which it is straightforward to find the answer, $2(a+b+35) \Rightarrow \frac{672}{5}$. Thus, $m+n = \frac{672}{5}\implies\boxed{677}$

 





\textbf{Solution 7}

We will bash with trigonometry.

 Firstly, by Pythagoras Theorem, $PQ=QR=RS=SP=25$. We observe that $[PQRS]=\frac{1}{2}\cdot30\cdot40=600$. Thus, if we drop an altitude from $P$ to $\overline{SR}$ to point $E$, it will have length $\frac{600}{25}=24$. In particular, $SE=7$ since we form a 7-24-25 triangle. 

 Now, $\sin\angle APS=\sin\angle SPB=\sin(\angle SPQ+\angle QPB)=\sin\angle SPQ\cos\angle QPB+\sin\angle QPB\cos\angle SPQ=\sin\angle PSR\cos\angle QPB-\sin\angle QPB\cos\angle PSR=\frac{24}{25}\cdot\frac{15}{25}-\frac{20}{25}\cdot\frac{7}{25}=\frac{44}{125}$. Thus, since $PS=25$, we get that $AS=\frac{44}{5}$. Now, by the Pythagorean Theorem, $AP=\frac{117}{5}$. 

 Using the same idea, $\cos\angle RSD=-\cos\angle RSA=-\cos(\angle RSP+\angle PSA)=\sin\angle RSP\sin\angle PSA-\cos\angle RSP\cos\angle PSA=\frac{24}{25}\cdot\frac{117}{125}-\frac{7}{25}\cdot\frac{44}{125}=\frac{4}{5}$. Thus, since $SR=20$. 

 Now, we can finish. We know $AB=\frac{117}{5}+15=\frac{192}{5}$. We also know $AD=\frac{44}{5}+20=\frac{144}{5}$. Thus, our perimeter is $\frac{672}{5}\implies\boxed{677}$

 





\textbf{Solution 8 (Pythagorean Theorem Bash)}


                    \begin{figure}[H]
                        \centering
                        \includegraphics[width=0.5\textwidth]{Solution/AIME_1991_I_12_Solution_Figure_1.png}
                    \end{figure}
                    Note that this is a modified version of the original diagram. By the Pythagorean theorem, the side length of the rhombus is $25$. Since a rhombus is a parallelogram, its diagonals bisect each other. Thus, $\triangle PBQ \cong \triangle RDS$ and $\triangle QCR \cong \triangle SAP$ by symmetry.

 Let $RC=x$ and $QC=y$. Applying the Pythagorean theorem to $\triangle QCR$, we obtain
 \[x^2+y^2=625.\]

 Next, draw a point $F$ such that $QF\perp AD$. Applying the Pythagorean theorem again gives
 \[(20-y)^2+(15+x)^2=1600.\]Expanding,
 \[y^2-40y+400+x^2+30x+225=1600.\]Substituting $x^2+y^2=625$ and combining constants yields
 \[30x-40y=350,\]which simplifies to
 \[3x=4y+35 \quad \Rightarrow \quad x=\frac{4y+35}{3}.\]

 Substituting this into $x^2+y^2=625$, we get
 \[\left(\frac{4y+35}{3}\right)^2+y^2=625.\]Multiplying through by $9$,
 \[25y^2+280y+1225=5625.\]Simplifying,
 \[5y^2+56y-880=0.\]Factoring,
 \[(5y-44)(y+20)=0.\]Since $y>0$, we have
 \[y=\frac{44}{5}.\]

 Substituting this into $x^2+y^2=625$,
 \[x^2+\frac{44^{2}}{25}=\frac{125^{2}}{25},\]so
 \[x^2=\frac{125^{2}-44^{2}}{25}=\frac{(125+44)(125-44)}{25}=\frac{169\cdot81}{25},\]giving
 \[x=\frac{117}{5}.\]

 Finally, the perimeter is
 \[2\left(20+15+\frac{44}{5}+\frac{117}{5}\right)=\frac{672}{5}.\]Since the problem asks for $m+n$, the answer is \(672+5=\boxed{677}.\)

 \textbf{Diagram} 
                    \begin{figure}[H]
                        \centering
                        \includegraphics[width=0.5\textwidth]{Solution/AIME_1991_I_12_Solution_Figure_1.png}
                    \end{figure}
                    




	\item (\textit{AIME I 1991}) A drawer contains a mixture of red socks and blue socks, at most 1991 in all. It so happens that, when two socks are selected randomly without replacement, there is a probability of exactly $\frac{1}{2}$ that both are red or both are blue. What is the largest possible number of red socks in the drawer that is consistent with this data? 


 
 \textbf{Answer: }990



\textbf{Solution 1}

Let $r$ and $b$ denote the number of red and blue socks, respectively. Also, let $t=r+b$. The probability $P$ that when two socks are drawn randomly, without replacement, both are red or both are blue is given by

 
 \[\frac{r(r-1)}{(r+b)(r+b-1)}+\frac{b(b-1)}{(r+b)(r+b-1)}=\frac{r(r-1)+(t-r)(t-r-1)}{t(t-1)}=\frac{1}{2}.\]

 Solving the resulting quadratic equation $r^{2}-rt+t(t-1)/4=0$, for $r$ in terms of $t$, one obtains that 

 
 \[r=\frac{t\pm\sqrt{t}}{2}\, .\]

 Now, since $r$ and $t$ are positive integers, it must be the case that $t=n^{2}$, with $n\in\mathbb{N}$. Hence, $r=n(n\pm 1)/2$ would correspond to the general solution. For the present case $t\leq 1991$, and so one easily finds that $n=44$ is the largest possible integer satisfying the problem conditions. 

 In summary, the solution is that the maximum number of red socks is $r=\boxed{990}$.

 





\textbf{Solution 2}

Let $r$ and $b$ denote the number of red and blue socks such that $r+b\le1991$. Then by complementary counting, the number of ways to get a red and a blue sock must be equal to $1-\frac12=\frac12=\frac{2rb}{(r+b)(r+b-1)}\implies4rb=(r+b)(r+b-1)$$=(r+b)^2-(r+b)\implies r^2+2rb+b^2-r-b=4rb\implies r^2-2rb+b^2$$=(r-b)^2=r+b$, so $r+b$ must be a perfect square $k^2$. Clearly, $r=\frac{k^2+k}2$, so the larger $k$, the larger $r$: $k^2=44^2$ is the largest perfect square below $1991$, and our answer is $\frac{44^2+44}2=\frac12\cdot44(44+1)=22\cdot45=11\cdot90=\boxed{990}$.

 


 





\textbf{Solution 3}

Let $r$ and $b$ denote the number of red and blue socks, respectively. In addition, let $t = r + b$, the total number of socks in the drawer.

 From the problem, it is clear that $\frac{r(r-1)}{t(t-1)} + \frac{b(b-1)}{t(t-1)} = \frac{1}{2}$

 Expanding, we get $\frac{r^2 + b^2 - r - b}{t^2 - t} = \frac{1}{2}$

 Substituting $t$ for $r + b$ and cross multiplying, we get $2r^2 + 2b^2 - 2r - 2b = r^2 + 2br + b^2 - r - b$

 Combining terms, we get $b^2 - 2br + r^2 - b - r = 0$

 To make this expression factorable, we add $2r$ to both sides, resulting in $(b - r)^2 - 1(b-r) = (b - r - 1)(b - r) = 2r$

 From this equation, we can test values for the expression $(b - r - 1)(b - r)$, which is the multiplication of two consecutive integers, until we find the highest value of $b$ or $r$ such that $b + r \leq 1991$. 

 By testing $(b - r - 1) = 43$ and $(b - r) = 44$, we get that $r = 43(22) = 946$ and $b = 990$. Testing values one integer higher, we get that $r = 990$ and $b = 1035$. Since $990 + 1035 = 2025$ is greater than $1991$, we conclude that $(946, 990)$ is our answer. 

 Since it doesn't matter whether the number of blue or red socks is $990$, we take the lower value for $r$, thus the maximum number of red socks is $r=\boxed{990}$.

 





\textbf{Solution 4}

As above, let $r$, $b$, and $t$ denote the number of red socks, the number of blue socks, and the total number of socks, respectively. We see that $\frac{r(r-1)}{t(t-1)}+\frac{b(b-1)}{t(t-1)}=\frac{1}{2}$, so $r^2+b^2-r-b=\frac{t(t-1)}{2}=r^2+b^2-t=\frac{t^2}{2}-\frac{t}{2}$. 

 Seeing that we can rewrite $r^2+b^2$ as $(r+b)^2-2rb$, and remembering that $r+b=t$, we have $\frac{t^2}{2}-\frac{t}{2}=t^2-2rb-t$, so $2rb=\frac{t^2}{2}-\frac{t}{2}$, which equals $r^2+b^2-t$.

 We now have $r^2+b^2-t=2rb$, so $r^2-2rb+b^2=t$ and $r-b=\pm\sqrt{t}$. Adding this to $r+b=t$, we have $2r=t\pm\sqrt{t}$. To maximize $r$, we must use the positive square root and maximize $t$. The largest possible value of $t$ is the largest perfect square less than 1991, which is $1936=44^2$. Therefore, $r=\frac{t+\sqrt{t}}{2}=\frac{1936+44}{2}=\boxed{990}$.

 Solution by Zeroman

 





\textbf{Solution 5}

Let $r$ be the number of socks that are red, and $t$ be the total number of socks. We get:

 $2(r(r-1)+(t-r)(t-r-1))=t(t-1)$Expanding the left hand side and the right hand side, we get:$4r^2-4rt+2t^2-2t = t^2-t$

 And, moving terms, we will get that:$4r^2-4rt+t^2 = t$

 We notice that the left side is a perfect square. $(2r-t)^2 = t$

 Thus $t$ is a perfect square. And, the higher $t$ is, the higher $r$ will be. So, we should set $t = 44^2 = 1936$

 And, we see, $2r-1936 = \pm44$ We will use the positive root, to get that $2r-1936 = 44$, $2r = 1980$, and $r = \boxed{990}$.

 - AlexLikeMath

 





\textbf{Solution 6}

Let $r$ and $b$ denote the red socks and blue socks, respectively. Thus the equation in question is:

 $\frac{r(r-1)+b(b-1)}{(r+b)(r+b-1)}=\frac{1}{2}$

 $\Rightarrow 2r^2-2r+2b^2-2b=r^2+2rb+b^2-r-b$

 $\Rightarrow r^2+b^2-r-b-2rb=0$

 $\Rightarrow (r-b)^2=r+b\le 1991$

 Because we wish to maximize $r$, we have $r\ge b$ and thus $r-b\le 44$ as both $r$ and $b$ must be integers and ${44}^2=1936$ is the largest square less than or equal to $1991$. We now know that the maximum difference between the number of socks is $44$. Now we return to an earlier equation:

 $r^2+b^2-r-b-2rb=0$

 $\Rightarrow r^2-(2b+1)r+(b^2-b)=0$

 Solving by the \href{https://artofproblemsolving.com/wiki/index.php?title=Quadratic\_formula}{Quadratic formula}, we have:

 $r=\frac{2b+1\pm\sqrt{4b^2+4b+1-4b^2+4b}}{2}$

 $\Rightarrow r=\frac{2b+1\pm\sqrt{8b+1}}{2}$

 $\Rightarrow r-b=\frac{1\pm\sqrt{8b+1}}{2}$

 $\Rightarrow 44\ge r-b=\frac{1\pm\sqrt{8b+1}}{2}$

 $\Rightarrow 44\ge\frac{1\pm\sqrt{8b+1}}{2}$

 $\Rightarrow \pm\sqrt{8b+1}\le 87$

 $\Rightarrow b\le 946$ (Here we use the positive sign to maximize $r$.)

 Thus the optimal case would be $b=946$ as $r$ increases only when $b$ increases. In addition, $\sqrt{8b+1}=87$ as we are using the equality case. We then plug back in:

 $r=\frac{2b+1\pm\sqrt{8b+1}}{2}$

 $\Rightarrow r=\frac{1893+87}{2}$ (using the established $\sqrt{8b+1}=87$)

 $\Rightarrow r=\frac{1980}{2}$

 $\Rightarrow r=\boxed{990}$

 




	\item (\textit{AIME I 1991}) A hexagon is inscribed in a circle. Five of the sides have length 81 and the sixth, denoted by $\overline{AB}$, has length 31. Find the sum of the lengths of the three diagonals that can be drawn from $A_{}^{}$.


 
 \textbf{Answer: }384



\textbf{Solution}


                    \begin{figure}[H]
                        \centering
                        \includegraphics[width=0.5\textwidth]{Solution/AIME_1991_I_14_Solution_Figure_1.png}
                    \end{figure}
                    Let $x=AC=BF$, $y=AD=BE$, and $z=AE=BD$.

 \href{https://artofproblemsolving.com/wiki/index.php?title=Ptolemy\%27s\_Theorem}{Ptolemy's Theorem} on $ABCD$ gives $81y+31\cdot 81=xz$, and Ptolemy on $ACDF$ gives $x\cdot z+81^2=y^2$.Subtracting these equations give $y^2-81y-112\cdot 81=0$, and from this $y=144$. Ptolemy on $ADEF$ gives $81y+81^2=z^2$, and from this $z=135$. Finally, plugging back into the first equation gives $x=105$, so $x+y+z=105+144+135=\boxed{384}$.

 





\textbf{Solution 2}

Let $\theta$ be the inscribed angle in each of the 5 sides of length 81, so $d \sin \theta = 81$. Since the inscribed angles sum to $\pi$, $d \sin 5\theta = d \sin (\pi - 5\theta) = 31$.

 Now consider the Chebyshev polynomials that put $\dfrac{\sin n\theta}{\sin \theta}$ in terms of $\cos \theta$:

 $\dfrac{\sin 2\theta}{\sin \theta} = 2 \cos \theta, \dfrac{\sin 3\theta}{\sin \theta} = 4 \cos^2 \theta - 1$

 $\dfrac{\sin 4\theta}{\sin \theta} = 8 \cos^3 \theta - 4 \cos \theta, \dfrac{\sin 5\theta}{\sin \theta} = 16 \cos^4 \theta - 12 \cos^2 \theta + 1$

 The sum of the diagonals is $d\sin 2\theta + d\sin 3\theta + d\sin 4\theta$, which becomes $(d \sin \theta)(8\cos^3 \theta + 4\cos^2 \theta - 2\cos \theta - 1)$, and we're given $16 \cos^4 \theta - 12 \cos^2 \theta + 1 = \dfrac{31}{81}$

 Solve for $\cos \theta$: $16 \cos^4 \theta - 12 \cos^2 \theta + \dfrac{9}{4} = \dfrac{9}{4} - \dfrac{50}{81}$

 $\left(4 \cos^2 \theta - \dfrac{3}{2}\right)^2 = \dfrac{729}{324} - \dfrac{200}{324} = \left(\dfrac{23}{18}\right)^2$

 $8 \cos^2 \theta - 3 = \dfrac{23}{9}$ or $-\dfrac{23}{9}$, so $8 \cos^2 \theta = \dfrac{50}{9}$ or $\dfrac{4}{9}$

 $\cos^2 \theta = \dfrac{25}{36}$ or $\dfrac{1}{18}$, which means $\cos \theta$ must be $\dfrac{5}{6}$ if $5 \theta < \pi$.

 Now $(d \sin \theta)\left(8\cos^3 \theta + 4\cos^2 \theta - 2\cos \theta - 1\right) = 81\left(8 \cdot \dfrac{125}{216} + 4 \cdot \dfrac{25}{36} - 2 \cdot \dfrac{5}{6} - 1\right)$

 $= 3 \left(8 \cdot \dfrac{125}{8} + 4 \cdot \dfrac{75}{4} - 2 \cdot \dfrac{45}{2} - 27\right) = 3(125 + 75 - 45 - 27) = \boxed{384}$

 





\textbf{Solution 3 (Risky)}

This solution refers to the diagram given in Solution 1.

 Applying Ptolemy’s Theorem to quadrilateral \(ADEF\), we obtain
 \[81y + 81^2 = z^2.\]

 Applying Ptolemy’s Theorem to quadrilateral \(ABCF\), we obtain
 \[31y + 81^2 = x^2.\]

 Since this is an AIME problem and the question asks for the sum of the lengths of the three diagonals, we may assume that all diagonal lengths are integers.

 Rewriting the first equation,
 \[81y = z^2 - 81^2 = (z+81)(z-81).\]

 Because both \(y\) and \(z\) are integers, \(z\) must be a multiple of \(9\). Additionally, we observe that \(y > z\), so we test the smallest such value of \(z\) satisfying this condition, namely \(z = 135\). Substituting gives
 \[y = \frac{135^2 - 81^2}{81} = 144.\]

 Substituting \(y = 144\) into the second equation,
 \[x^2 = 31(144) + 81^2 = 11025,\]so \(x = 105\).

 Finally, summing the three diagonal lengths, \(105 + 135 + 144 = \boxed{384}.\)

 




	\item (\textit{AIME I 1991}) For positive integer $n_{}^{}$, define $S_n^{}$ to be the minimum value of the sum


 $\sum_{k=1}^n \sqrt{(2k-1)^2+a_k^2},$where $a_1,a_2,\ldots,a_n^{}$ are positive real numbers whose sum is 17. There is a unique positive integer $n^{}_{}$ for which $S_n^{}$ is also an integer. Find this $n^{}_{}$.


 
 \textbf{Answer: }012



\textbf{Solution 1 (Geometric Interpretation)}

Consider $n$ right triangles joined at their vertices, with bases $a_1,a_2,\ldots,a_n$ and heights $1,3,\ldots, 2n - 1$. The sum of their hypotenuses is the value of $S_n$. The minimum value of $S_n$, then, is the length of the straight line connecting the bottom vertex of the first right triangle and the top vertex of the last right triangle, so
 \[S_n \ge \sqrt {\left(\sum_{k = 1}^n (2k - 1)\right)^2 + \left(\sum_{k = 1}^n a_k\right)^2}.\]Since the sum of the first $n$ odd integers is $n^2$ and the sum of $a_1,a_2,\ldots,a_n$ is 17, we get
 \[S_n \ge \sqrt {17^2 + n^4}.\]If this is an integer, we can write $17^2 + n^4 = m^2$, for an integer $m$. Thus, $(m - n^2)(m + n^2) = 289\cdot 1 = 17\cdot 17 = 1\cdot 289.$ The only possible value, then, for $m$ is $145$, in which case $n^2 = 144$, and $n = \boxed {012}$.

 





\textbf{Solution 2}

The inequality 
 \[S_n \ge \sqrt {\left(\sum_{k = 1}^n (2k - 1)\right)^2 + \left(\sum_{k = 1}^n a_k\right)^2}.\]is a direct result of the \href{https://artofproblemsolving.com/wiki/index.php?title=Minkowski\_Inequality}{Minkowski Inequality}. Continue as above.

 





\textbf{Solution 3}

Let $a_{i} = (2i - 1) \tan{\theta_{i}}$ for $1 \le i \le n$ and $0 \le \theta_{i} < \frac {\pi}{2}$. We then have that
 \[S_{n} = \sum_{k = 1}^{n}\sqrt {(2k - 1)^{2} + a_{k}^{2}} = \sum_{k = 1}^{n}(2k - 1) \sec{\theta_{k}}\]Note that that $S_{n} + 17 = \sum_{k = 1}^{n}(2k - 1)(\sec{\theta_{k}} + \tan{\theta_{k}})$.Note that for any angle $\theta$, it is true that $\sec{\theta} + \tan{\theta}$ and $\sec{\theta} - \tan{\theta}$ are reciprocals. We thus have that $S_{n} - 17 = \sum_{k = 1}^{n}(2k - 1)(\sec{\theta_{k}} - \tan{\theta_{k}}) = \sum_{k = 1}^{n}\frac {2k - 1}{\sec{\theta_{k}} + \tan{\theta_{k}}}$. By the AM-HM inequality on these $n^{2}$ values, we have that:
 \[\frac {S_{n} + 17}{n^{2}}\ge \frac {n^{2}}{S_{n} - 17}\rightarrow S_{n}^{2}\ge 289 + n^{4}\]This is thus the minimum value, with equality when all the tangents are equal. The only value for which $\sqrt {289 + n^{4}}$ is an integer is $n = 12$ (see above solutions for details).

 





\textbf{Solution 4 (Vector)}

Observe the sum and apply it to the Cartesian plane, we can see that it is basically calculating distance between the origin and two points. Therefore we can assume every $\sqrt{(2k-1)^2+a_k^2}$ to be a vector from the origin to point $(2k-1, a_k)$. Now we can do sum inside the vector so we get $((1+3+5+\ldots+2k-1,   a_1+a_2+a_3+\ldots+a_k)$, with $a_1+\ldots+a_k$ to be $17$ and $1+3+\ldots+2k-1$ to be $k^2$. Then we calculate the length of the vector to be $\sqrt{k^4+289}$. Since the sum needs to be an integer, we assume $k^4+289$ equals to $s^2$. Applying difference between squares, we get that $k^2-s=289$ and $k^2+s=-1$. Therfore $2k^2$ is $288$ and $k=n=12$ to be the final answer.

 JINZHENQIAN

 




\end{enumerate}